\chapter{Laryngoscopy}

\section*{Definition and Overview}
Laryngoscopy is a medical procedure used to examine the larynx, which houses the vocal cords, as well as the adjacent throat structures. It is a fundamental tool in otorhinolaryngology (ear, nose, and throat medicine), speech pathology, and anaesthesia, allowing clinicians to directly or indirectly visualise the upper airway. Laryngoscopy can be performed for diagnostic purposes, such as investigating chronic hoarseness, difficulty swallowing, or throat pain, and for therapeutic interventions, including the removal of foreign bodies, biopsy of suspicious lesions, or facilitation of endotracheal intubation. The procedure is broadly classified into indirect laryngoscopy (using a mirror or a flexible nasendoscope in an outpatient setting) and direct laryngoscopy (using a rigid laryngoscope, typically under general anaesthesia or in emergency airway management).

\section*{Epidemiology}
As a diagnostic and therapeutic procedure, laryngoscopy is performed very frequently across diverse clinical settings worldwide. In Australia, thousands of flexible nasolaryngoscopies are conducted annually in outpatient clinics, particularly among adults presenting with voice disorders, dysphagia, or suspected head and neck malignancies. Vocal cord lesions and conditions requiring laryngoscopy, such as nodules, polyps, or laryngeal cancer, show varied demographic distributions; for instance, vocal nodules are more prevalent in professional voice users and school-aged children, whereas laryngeal squamous cell carcinoma occurs predominantly in males over the age of 50 with a history of heavy smoking and alcohol consumption. In emergency and critical care departments, direct laryngoscopy is a daily necessity for securing airways in critically ill or injured patients.

\section*{Aetiology and Pathophysiology}
Laryngoscopy is indicated when there is a suspected pathological process affecting the larynx, pharynx, or upper airway. The underlying mechanisms and clinical indications that necessitate this procedure include:
\begin{itemize}
    \item \textbf{Vocal fold pathology:} organic lesions such as vocal nodules, cysts, polyps, contact granulomas, or papillomas that disrupt the mucosal wave of the vocal folds, leading to dysphonia.
    \item \textbf{Malignancy screening:} suspected neoplastic changes, particularly squamous cell carcinoma of the larynx, which are often triggered by chronic mucosal irritation from tobacco smoke, alcohol, and human papillomavirus (HPV) infection.
    \item \textbf{Inflammatory conditions:} acute or chronic laryngitis, laryngopharyngeal reflux (LPR), or autoimmune conditions causing laryngeal inflammation and mucosal thickening.
    \item \textbf{Neurological dysfunction:} vocal fold paralysis or paresis resulting from recurrent laryngeal nerve damage (often post-thyroidectomy or due to mediastinal pathology) or central neurological disorders like stroke or Parkinson's disease.
    \item \textbf{Airway obstruction:} physical compromise of the upper airway from epiglottitis, laryngeal oedema (e.g., anaphylaxis), foreign body inhalation, or external trauma to the neck.
    \item \textbf{Intubation requirements:} the physiological necessity to establish a secure airway for mechanical ventilation during surgery or critical illness.
\end{itemize}

\section*{Clinical Features}
Laryngoscopy is initiated in response to specific clinical signs and symptoms indicating laryngeal or pharyngeal dysfunction. These clinical features include:
\begin{itemize}
    \item \textbf{Dysphonia:} persistent or progressive hoarseness, breathiness, or loss of voice lasting longer than three weeks.
    \item \textbf{Stridor:} a high-pitched, noisy breathing sound indicating upper airway obstruction, which requires urgent evaluation.
    \item \textbf{Dysphagia:} difficulty or pain when swallowing (odynophagia), particularly when accompanied by weight loss or ear pain (referred otalgia).
    \item \textbf{Globus sensation:} a persistent feeling of a lump or foreign body in the throat, unrelated to swallowing.
    \item \textbf{Chronic cough:} an unexplained cough lasting for months that has not responded to standard pulmonary therapies.
    \item \textbf{Haemoptysis:} coughing up blood or blood-stained mucus, which may originate from a laryngeal or pharyngeal lesion.
\end{itemize}

\section*{Differential Diagnosis}
While laryngoscopy is a diagnostic tool, the clinician must consider alternative or complementary diagnostic modalities when evaluating upper airway symptoms:
\begin{itemize}
    \item \textbf{Bronchoscopy:} indicated when symptoms like chronic cough or haemoptysis are suspected to originate from the trachea, bronchi, or lungs rather than the larynx.
    \item \textbf{Oesophagogastroduodenoscopy (OGD):} preferred when dysphagia is suspected to be secondary to distal oesophageal pathology, such as strictures, achalasia, or gastro-oesophageal reflux disease (GORD).
    \item \textbf{Barium swallow study:} a video-fluoroscopic examination useful for evaluating the mechanical and functional phases of swallowing when structural laryngeal pathology is not primary.
    \item \textbf{High-resolution CT or MRI of the neck:} indicated for staging suspected laryngeal tumours or evaluating deep neck space infections that cannot be fully visualised endoscopically.
    \item \textbf{Videostroboscopy:} a specialised form of laryngoscopy utilising pulsed light to assess the mucosal wave and vocal fold vibration in patients with complex voice disorders.
\end{itemize}

\section*{When to Seek Medical Care}
Patients experiencing persistent voice changes for more than three weeks, unexplained throat pain, or difficulty swallowing should consult a general practitioner or an ear, nose, and throat (ENT) specialist for evaluation, which may include a laryngoscopy.

\textbf{Call 000 (or local emergency services) for:}
\begin{itemize}
    \item Severe difficulty breathing, gasping for air, or stridor (noisy, high-pitched breathing)
    \item Sudden swelling of the lips, tongue, or throat, especially with a history of allergic exposure (anaphylaxis)
    \item Inability to swallow saliva, leading to drooling
    \item Cyanosis (bluish discolouration of the lips, face, or skin) indicating oxygen deprivation
    \item Neck trauma accompanied by voice change, breathing difficulties, or rapid neck swelling
\end{itemize}

\section*{Diagnosis and Investigations}
The decision to perform a laryngoscopy is made following a thorough initial evaluation, and the procedure itself serves as a key diagnostic investigation. The diagnostic pathway includes:
\begin{itemize}
    \item \textbf{Clinical history:} a detailed history focusing on the duration and onset of symptoms, vocal hygiene, occupational voice use, smoking and alcohol history, and exposure to reflux triggers.
    \item \textbf{Physical examination:} palpation of the neck for lymphadenopathy, thyroid nodules, or laryngeal crepitus, alongside an examination of the oral cavity and oropharynx.
    \item \textbf{Indirect laryngoscopy:} using a small laryngeal mirror positioned in the back of the throat to reflect light and provide a basic view of the vocal folds.
    \item \textbf{Flexible nasolaryngoscopy:} a key diagnostic test where a thin, flexible endoscope is passed through the nasal cavity to provide a detailed, high-definition view of the larynx during speech and swallowing.
    \item \textbf{Direct rigid laryngoscopy:} performed in an operating theatre under general anaesthesia, using a rigid hollow metal scope to allow direct visualisation, biopsy of lesions, or microsurgical interventions.
    \item \textbf{Laryngeal biopsy:} the extraction of tissue samples during direct laryngoscopy for histopathological analysis, essential for confirming or ruling out malignancy.
\end{itemize}

\section*{Management}
The management of laryngoscopy involves proper patient preparation, choice of appropriate technique, and post-procedure care:
\begin{itemize}
    \item \textbf{Procedural preparation:} obtaining informed consent, explaining the procedure, and applying a topical decongestant and local anaesthetic spray (such as lidocaine) to the nasal cavity and pharynx to minimise discomfort and suppress the gag reflex.
    \item \textbf{Anaesthetic management:} direct rigid laryngoscopy requires general anaesthesia with endotracheal intubation or jet ventilation, requiring coordination with an anaesthetist to maintain airway safety.
    \item \textbf{Post-procedure monitoring:} observing the patient in a recovery area until local anaesthetic effects wear off, ensuring they do not eat or drink for at least one hour to prevent aspiration.
    \item \textbf{Voice rest:} advising short-term voice rest or voice conservation if a biopsy or vocal cord surgery was performed.
    \item \textbf{Pharmacological adjuncts:} prescribing simple analgesics for throat soreness, and proton pump inhibitors or lifestyle modifications if findings indicate laryngopharyngeal reflux.
    \item \textbf{Speech pathology referral:} referring patients with functional voice disorders, nodules, or muscle tension dysphonia for targeted voice therapy.
\end{itemize}

\section*{Complications}
While laryngoscopy is generally safe, potential complications can arise depending on the type of procedure performed:
\begin{itemize}
    \item Mucosal trauma, including minor bleeding or epistaxis from the nose or throat
    \item Dental injury, such as chipped teeth, which is a specific risk during rigid direct laryngoscopy
    \item Laryngospasm or bronchospasm, particularly in patients with hyperreactive airways (such as asthma or active respiratory infection)
    \item Aspiration of saliva or stomach contents if the patient attempts to eat or drink before the topical anaesthetic has worn off
    \item Adverse reactions to local or general anaesthetic agents, including allergic reactions or respiratory depression
    \item Temporary hoarseness or exacerbation of throat pain following a biopsy or instrumentation
\end{itemize}

\section*{Prognosis and Outcomes}
The prognosis and outcomes of laryngoscopy are excellent, as it is a highly reliable diagnostic tool with a low rate of serious complications. Flexible nasolaryngoscopy is well-tolerated by the vast majority of patients and provides immediate visual results that guide management. For patients undergoing direct rigid laryngoscopy and biopsy under general anaesthesia, recovery is typically rapid, with most patients returning home on the same day. The long-term prognosis depends entirely on the underlying pathology diagnosed during the procedure, such as whether a laryngeal lesion is benign (e.g., vocal nodules, which respond well to therapy) or malignant (which requires staging and multidisciplinary oncology care).

\section*{Prevention and Self-Care}
Patients can take several steps to prepare for laryngoscopy and care for themselves after the procedure:
\begin{itemize}
    \item Fast from food and fluids for the recommended period before the procedure, especially if general anaesthesia or deep sedation is planned
    \item Avoid eating or drinking for at least one hour after outpatient flexible laryngoscopy until normal throat sensation returns
    \item Practice strict voice rest (no whispering or shouting) if advised by the ENT specialist following a biopsy or surgery
    \item Hydrate well by sipping cool water once the anaesthetic has worn off, and use saline gargles to soothe throat soreness
    \item Avoid smoking and exposure to second-hand smoke, which irritate the mucosal lining of the throat and larynx
    \item Seek immediate medical attention if you experience breathing difficulties, coughing up large amounts of blood, or worsening neck swelling post-procedure
\end{itemize}

\section*{Sources and Further Reading}
The following reputable organisations provide further information on laryngoscopy and related throat procedures:
\begin{itemize}
    \item Healthdirect Australia. \textit{Laryngoscopy procedure and patient guide.} https://www.healthdirect.gov.au/
    \item Mayo Clinic. \textit{Laryngoscopy information, preparation, and risks.} https://www.mayoclinic.org/
    \item NHS. \textit{Laryngoscopy: what to expect during the test.} https://www.nhs.uk/
    \item Australian Society of Otolaryngology Head and Neck Surgery. \textit{Patient resources on throat and airway examinations.} https://asohns.org.au/
    \item MedlinePlus. \textit{Laryngoscopy patient instructions.} https://medlineplus.gov/
\end{itemize}
